% part: intuitionistic-logic
% chapter: introduction
% section: natural-deduction

\documentclass[../../../include/open-logic-section]{subfiles}

\begin{document}

\olfileid{int}{int}{ntd}

\olsection{فطری استنباط}

$\FalseCl$ قاعدیاں توں بغیر فطری استنباط وجدانیت پسند منطق دا اک
معیاری !!{derivation} نظام اے۔ ایتھے اسیں قاعدے فیر دیندے آں تے بی
ایچ کے تعبیر نال اوہناں دی وجھ دسّنے آں۔ ہر صورت وچ اسیں قاعدے نوں
ایویں سمجھ سکدے آں کہ جے مقدمیاں دیاں تعمیراں نیں، تاں نتیجے دی
تعمیر وی اے۔

فطری استنباط دیاں !!{derivation}s وچ ہجے قائم مفروضے ہوندے نیں۔
سو ایسے !!a{derivation} نوں، مثلاً $!A$ دی !!{undischarged}
مفروضیاں~$\Gamma$ توں، اک فنکشن سمجھنا چاہیدا اے جیہڑا سارے
$!B \in \Gamma$ دیاں تعمیراں نوں~$!A$ دی تعمیر وچ بدلدا اے۔
جے !!a{derivation}~$!A$ دی کسے وی !!{undischarged} مفروضے توں بغیر
اے، تاں بی ایچ کے تعبیر دے معنی وچ~$!A$ دی تعمیر اے۔ گل سمجھاؤن
لئی جدوں لوڑ نہیں ہووے گی، اسیں~$\Gamma$ چھڈ دیاں گے۔

مفروضہ~$!A$ اپنے آپ وچ !!a{derivation} اے،~$!A$ دی
!!{undischarged} مفروضے~$!A$ توں۔ ایہ بی ایچ کے تعبیر نال میل
کھاندا اے۔ تعمیراں اُتے شناختی فنکشن~$!A$ دی ہر تعمیر نوں~$!A$
دی تعمیر وچ بدلدا اے۔

\subsection{عطف}

\begin{defish}
\AxiomC{$!A$}
\AxiomC{$!B$}
\RightLabel{\Intro{\land}}
\BinaryInfC{$!A \land !B$}
\DisplayProof
\hfill
\begin{tabular}{l}
\AxiomC{$!A \land !B$}
\RightLabel{\Elim{\land}}
\UnaryInfC{$!A$}
\DisplayProof
\\[3ex]
\AxiomC{$!A \land !B$}
\RightLabel{\Elim{\land}}
\UnaryInfC{$!B$}
\DisplayProof
\end{tabular}
\end{defish}

منّ لو ساڈے کول تعمیراں $N_1$، $N_2$ نیں، واری واری $!A_1$ تے
$!A_2$ دیاں۔ تاں ساڈے کول $!A_1 \land !A_2$ دی تعمیر وی اے،
یعنی جوڑی $\tuple{N_1, N_2}$۔

بی ایچ کے تعبیر مطابق $!A_1 \land !A_1$ دی تعمیر اک جوڑی
$\tuple{N_1, N_2}$ اے۔ سو منّ لو ساڈے کول ایہو جیہی جوڑی اے۔
تاں ساڈے کول عطف دے ہر جز دی تعمیر وی اے۔ $N_1$،~$!A_1$
دی تعمیر اے تے $N_2$، $!A_2$ دی تعمیر اے۔

\subsection{شرطیہ}

\begin{defish}
  \AxiomC{$\Discharge{!A}{u}$}
  \DeduceC{$!B$}
  \DischargeRule{$\Intro{\lif}$}{u}
  \UnaryInfC{$!A \lif !B$}
  \DisplayProof
\hfill
  \AxiomC{$!A \lif !B$}
  \AxiomC{$!A$}
  \RightLabel{$\Elim{\lif}$}
  \BinaryInfC{$!B$}
  \DisplayProof
\end{defish}

جے ساڈے کول !!a{derivation}~$!B$ دی !!{undischarged}
مفروضے~$!A$ توں اے، تاں ایسا فنکشن~$f$ موجود اے جیہڑا
~$!A$ دیاں تعمیراں نوں~$!B$ دیاں تعمیراں وچ بدلدا اے۔ اوہی فنکشن
$!A \lif !B$ دی تعمیر اے۔ سو جے $\Intro\lif$ دے مقدمے دی
تعمیر~$!A$ دی تعمیر اُتے منحصر اے، تاں نتیجے $!A
\lif !B$ دی تعمیر اے۔

دوجے پاسے منّ لو تعمیر~$N$، $!A$ دی، تے تعمیر $f$،
$!A \lif !B$ دی، موجود نیں۔ $!A \lif !B$ دی تعمیر ایسا فنکشن
اے جیہڑا $!A$ دیاں تعمیراں نوں~$!B$ دیاں تعمیراں وچ بدلدا اے۔
سو $f(N)$،~$!B$ دی تعمیر اے، یعنی $\Elim\lif$ دے نتیجے دی تعمیر اے۔

\subsection{فصل}

\begin{defish}
\begin{tabular}{l}
\AxiomC{$!A$}
\RightLabel{\Intro{\lor}}
\UnaryInfC{$!A \lor !B$}
\DisplayProof
\\[3ex]
\AxiomC{$!B$}
\RightLabel{\Intro{\lor}}
\UnaryInfC{$!A \lor !B$}
\DisplayProof
\end{tabular}
\hfill
\AxiomC{$!A \lor !B$}
\AxiomC{$\Discharge{!A}{n}$}
\DeduceC{$!C$}
\AxiomC{$\Discharge{!B}{n}$}
\DeduceC{$!C$}
\DischargeRule{\Elim{\lor}}{n}
\TrinaryInfC{$!C$}
\DisplayProof
\end{defish}

جے ساڈے کول تعمیر~$N_i$،~$!A_i$ دی، اے، تاں اینوں
تعمیر~$\tuple{i, N_i}$،~$!A_1 \lor !A_2$ دی، وچ بدل سکدے آں۔
دوجے پاسے منّ لو ساڈے کول~$!A_1 \lor !A_2$ دی تعمیر اے، یعنی
اک جوڑی $\tuple{i, N_i}$ جتھے $N_i$،~$!A_i$ دی تعمیر اے۔
نال فنکشن $f_1$، $f_2$ وی نیں جیہڑے واری واری $!A_1$، $!A_2$
دیاں تعمیراں نوں~$!C$ دیاں تعمیراں وچ بدلدے نیں۔ تاں $f_i(N_i)$،
~$!C$ دی تعمیر اے، یعنی~$\Elim\lor$ دے نتیجے دی۔

\subsection{محال}

\begin{defish}
  \AxiomC{$\lfalse$}
  \RightLabel{\FalseInt}
  \UnaryInfC{$!A$}
  \DisplayProof
\end{defish}

جے ساڈے کول !!a{derivation}~$\lfalse$ دی !!{undischarged}
مفروضیاں~$!B_1$، \dots، $!B_n$ توں اے، تاں اک فنکشن~$f(M_1,
\dots, M_n)$ اے جیہڑا $!B_1$، \dots، $!B_n$ دیاں تعمیراں نوں
~$\lfalse$ دی تعمیر وچ بدلدا اے۔ $\lfalse$ دی کوئی تعمیر نہیں،
ایس لئی ساریاں $!B_1$، \dots، $!B_n$ دیاں تعمیراں وی اکٹھیاں
نہیں ہو سکدیاں۔ سو $f$ دی ایہ خاصیت وی اے کہ \emph{جے} $M_1$، \dots،
$M_n$ واری واری $!B_1$، \dots، $!B_n$ دیاں تعمیراں نیں،
\emph{تاں} $f(M_1, \dots, M_n)$،~$!A$ دی تعمیر اے۔

\subsection{$\lnot$ دے قاعدے}

$\lnot !A$ نوں $!A \lif \lfalse$ دے طور اُتے متعین کیتا گیا اے،
سو عین لحاظ نال~$\lnot$ لئی قاعدیاں دی لوڑ نہیں۔ جے لوڑ ہوندی،
تاں اوہ ایویں ہوندے:

\begin{defish}
\AxiomC{$\Discharge{!A}{n}$}
\noLine
\DeduceC{$\lfalse$}
\DischargeRule{\Intro{\lnot}}{n}
\UnaryInfC{$\lnot !A$}
\DisplayProof
\hfill
\AxiomC{$\lnot !A$}
\AxiomC{$!A$}
\RightLabel{\Elim{\lnot}}
\BinaryInfC{$\lfalse$}
\DisplayProof
\end{defish}

\subsection{\usetoken{P}{derivation} دیاں مثالاں}

\begin{enumerate}
\item $\Proves !A \lif (\lnot !A \lif \lfalse)$, یعنی $\Proves !A
  \lif ((!A \lif \lfalse) \lif \lfalse)$
  \begin{prooftree}
    \AxiomC{$\Discharge{!A}{2}$}
    \AxiomC{$\Discharge{!A \lif \lfalse}{1}$}
    \RightLabel{\Elim\lif}
    \BinaryInfC{$\lfalse$}
    \DischargeRule{\Intro\lif}{1}
    \UnaryInfC{$(!A \lif \lfalse)\lif \lfalse$}
    \DischargeRule{\Intro\lif}{2}
    \UnaryInfC{$!A \lif (!A \lif \lfalse) \lif \lfalse$}
  \end{prooftree}

\item $\Proves ((!A \land !B) \lif !C) \lif (!A \lif (!B \lif !C))$
  \begin{prooftree}
    \AxiomC{$\Discharge{(!A \land !B) \lif !C}{3}$}
    \AxiomC{$\Discharge{!A}{2}$}
    \AxiomC{$\Discharge{!B}{1}$}
    \RightLabel{\Intro\land}
    \BinaryInfC{$!A \land !B$}
    \RightLabel{\Elim\lif}
    \BinaryInfC{$!C$}
    \DischargeRule{\Intro\lif}{1}
    \UnaryInfC{$!B \lif !C$}
    \DischargeRule{\Intro\lif}{2}
    \UnaryInfC{$!A \lif (!B \lif !C)$}
    \DischargeRule{\Intro\lif}{3}
    \UnaryInfC{$((!A \land !B) \lif !C) \lif (!A \lif (!B \lif !C))$}
  \end{prooftree}

\item $\Proves \lnot(!A \land \lnot !A)$, یعنی $\Proves (!A \land (!A
  \lif \lfalse))\lif \lfalse$
  \begin{prooftree}
    \AxiomC{$\Discharge{!A \land (!A \lif \lfalse)}{1}$}
    \RightLabel{\Elim\land}
    \UnaryInfC{$!A \lif \lfalse$}
    \AxiomC{$\Discharge{!A \land (!A \lif \lfalse)}{1}$}
    \RightLabel{\Elim\land}
    \UnaryInfC{$!A$}
    \RightLabel{\Elim\lif}
    \BinaryInfC{$\lfalse$}
    \DischargeRule{\Intro\lif}{1}
    \UnaryInfC{$(!A \land (!A \lif \lfalse))\lif \lfalse$}
  \end{prooftree}

\item $\Proves \lnot\lnot(!A \lor \lnot !A)$, یعنی $\Proves ((!A \lor
  (!A \lif \lfalse))\lif \lfalse)\lif \lfalse$
  \begin{prooftree}\hskip-3em
    \AxiomC{$\Discharge{(!A \lor (!A \lif \lfalse))\lif \lfalse}{2}$}
    \AxiomC{$\Discharge{(!A \lor (!A \lif \lfalse))\lif \lfalse}{2}$}
    \AxiomC{$\Discharge{!A}{1}$}
    \RightLabel{\Intro\lor}
    \UnaryInfC{$!A \lor (!A \lif \lfalse)$}
    \RightLabel{\Elim\lif}
    \BinaryInfC{$\lfalse$}
    \DischargeRule{\Intro\lif}{1}
    \UnaryInfC{$!A \lif \lfalse$}
    \RightLabel{\Intro\lor}
    \UnaryInfC{$!A \lor (!A \lif \lfalse)$}
    \RightLabel{\Elim\lif}
    \insertBetweenHyps{\hskip -4em}
    \BinaryInfC{$\lfalse$}
    \DischargeRule{\Intro\lif}{2}
    \UnaryInfC{$((!A \lor (!A \lif \lfalse))\lif \lfalse)\lif \lfalse$}
  \end{prooftree}
\end{enumerate}

\begin{prop}
  جے وجدانیت پسند منطق دے فطری استنباط وچ $\Gamma \Proves !A$،
  تاں کلاسیکی منطق وچ $\Gamma \Proves !A$۔ خاص کر کے جے
  $!A$ وجدانیت پسند قضیہ اے، تاں کلاسیکی قضیہ وی اے۔
\end{prop}

\begin{proof}
  فطری استنباط دا ہر قاعدہ کلاسیکی فطری استنباط دا وی قاعدہ اے۔
  سو وجدانیت پسند منطق دی ہر !!{derivation} کلاسیکی منطق وچ وی
  !!a{derivation} اے۔
\end{proof}

\begin{prob}
  وجدانیت پسند منطق وچ ہِیٹھلے !!{formula}s دیاں !!{derivation}s دیو:
  \begin{enumerate}
    \item $(\lnot !A \lor !B) \lif (!A \lif !B)$
    \item $\lnot\lnot\lnot !A \lif \lnot !A$
    \item $\lnot\lnot (!A \land !B) \liff (\lnot\lnot !A\land \lnot \lnot !B)$
    \item $\lnot (!A \lor !B) \liff (\lnot !A \land \lnot !B)$
    \item $(\lnot !A \lor \lnot !B) \lif \lnot (!A \land !B)$
    \item $\lnot\lnot ( !A \land !B) \lif  (\lnot\lnot !A \lor
    \lnot\lnot !B)$
  \end{enumerate}
\end{prob}

\end{document}
